\subsection{First Letter}

\begin{quotex}
Among the 19,000 pieces of correspondence found in \textbf{Carl Schmitt}'s personal library, there were eight letters from Julius Evola over a period of several years. There were none found in the opposite direction. From the letters, it is obvious that Evola was very interested in Schmitt's book on \textbf{Donoso Cortes}, whom they both regarded quite highly. As a public service, we have made available Donoso's book, \textit{Essays on Catholicism, Liberalism, and Socialism}\footnote{\url{https://www.gornahoor.net/library/CortesEssays.pdf}}. Since it is hardly likely that Evola was interested in Donoso's take on Catholicism, we can assume the other essays are more pertinent. At the very least, they are a counter-balance to the absurd meme floated by the Nouvelle Droite regarding the intellectual source of liberalism. More likely, Evola and Schmitt were impressed by Donoso's defense of authoritarianism.

The letters were written in German, using the polite form of ``you''. The translations here derive from an Italian translation of the original. Evola's first letter to Schmitt follows:
\end{quotex}


15 December 1951

Dear Professor!

I owe your address to Dr. Mohler; therefore, I am able to take the initiative---which I gave been thinking about for some time --- \emph{après le déluge}. I have often requested news about you:

primarily, our common friend, Prince Rohan, has assured me that you have at least physically moved beyond the period of the fall. Subsequently, I learned of your controversial ``reappearance'' and new works. Regarding this, I thank you very much for the booklet ``Recht und Raum'', that I received hear in Rome.

As far as my own person, things have not always gone the best for me: a war wound prevents me from walking and I can remain seated only for a few hours a day. I returned to Italy in 1948, but I've had to stay in nursing homes, so that I returned to my former home in Rome only last May. Subsequently a strange thing happened: I was ... arrested. At the margins of so-called ``neofascism'' (MSI --- \textbf{Movimento Sociale Italiano}) groups were formed that committed some foolish acts (bombings). Since my writings were read in those circles with great frequency and my person was rather highly regarded as a ``spiritual father'', they wanted to attribute to me the responsibility of a movement of which I knew almost nothing, and they accused me of the alleged defense of ``fascist'' ideas (``apologist for Fascism''). The story ended last month with my complete innocence; the only consequences were free publicity in my favour and a bad impression for the imposing, farsighted state police.

Apart from this, after my return I took up again unchanged my old activity in the teaching and politico-cultural fields (conservative revolutionary, as Dr. Mohler would say). Nevertheless, the situation here is not very easy and not only because of the Christian-social democracy, but also for the heavy legacy of the so-called ``second Fascism'' (republican and ``social'') of the neo-Fascists.

A short time ago Revolt against the Modern World was published in a new and expanded edition. The same for my works of a spiritual and ``esoteric'' character that had gone out of print and in the meantime other translations were published in English (The Doctrine of Awakening). It is a way of trifling the time.

After having briefly told you about me, I would be happy to know something of you and your projects, since I would give great value to remaining in contact with you. I would also have to ask you one thing: could you possibly procure for me a copy of your new writing on Donoso Cortes? I myself, in fact, am interested in this author and intend to deal with him in an essay --- Menschen und Trummer --- on which I am currently working.

Well, I offer you my best wishes for the upcoming Holiday.

I remain with old friendship.

P.S. Did something happen to your old house in Dahlem?

I will send you a pamphlet, Orientamenti \emph{Orientations}, which constituted the principle \emph{corpus delicti} body of the crime, evidence of my trial.


\flrightit{Posted on 2012-06-25 by Cologero}

